\section{Comparison of our quantum Boltzmann methods to others}
In this section we will give a brief overview of how our quantum Boltzmann methods compare to the other main competitive methods out there. This section is meant to give insight into why certain design choices were made and how it improves the algorithm. In the appendix we furthermore provide a detailed complexity analysis of our CQBM compared to other implementations. 

\subsection{Collisionless Quantum Boltzmann methods - comparison}
Our approach furthermore improves over the original QBM \cite{Todorova2020} by Todorova and Steijl in several key manners. First of all our quantum specular reflection step, even though seemingly more complicated, performs the correct reflection behavior in all corner cases. This is not the case for the original QBM, which shows incorrect reflection behavior for particles hitting the corners of an obstacle on a non-axis-aligned trajectory. On top of that, our specular reflection approach also ensures that the particles are set back into the flow domain before the start of a new timestep if they virtually traveled into an object. This also avoids incorrect behavior present in the original QBM formulation.
Furthermore our approach for identifying the internal grid points of obstacles at which particles need to be reflected is more efficient than the method used by Todorova and Steijl. This leads to an additional polynomial speed-up.

We finally show that our quantum implementation of the streaming step outperforms the methods of \cite{Todorova2020, Budinski2020, Budinski2021} in terms of the the amount of CNOT gates required to implement them. We intentionally choose for the latter as performance measure as they are either available as native gates or can be efficiently emulated. Hence, complexity estimates in terms of CNOT gates give a more realistic picture of the costs on a real quantum computer, whereas complexity estimates in terms of multi-controlled gates conceal the costs for decomposing multi-controlled gates into native ones.

Last but not least our method outperforms the current state-of-the-art approaches in the encoding of the particles' velocity. We propose a novel encoding of the velocity vector, due to which velocities can be flipped with a single NOT operation performed on a single qubit, whereas before $n$ NOT operations were required to flip the direction of the velocity encoded using $n$ qubits \cite{Todorova2020}.

\subsection{Full quantum Boltzmann method - space time encoding compared to other methods}
There are other methods that have implemented both a streaming and a collision step for the Boltzmann equation. In this section we will compare our method to what we will refer to as 'stop-and-go' methods and the Carleman linearization method, which both allow for solving the full Boltzmann equation on a quantum computer. Specifically for the 'stop-and-go' method we will consider the method of Budinski described in \cite{Budinski2021} and for the Carleman linearization approach we will consider the paper \cite{Sanavio2023}. \\

We will first consider the method \cite{Budinski2021}, which requires measurement and reinitialization at the end of each time step. Both the measurement and reinitialization procedure are exponentially expensive for in the total number of qubits \cite{Shende2004}. This implies that any potential quantum advantage is lost when performing multiple timesteps. Next to that the initalization procedure itself might already be exponentially expensive in the first timestep, as the function that is encoded is in no way trivial to encode. Therefore even when performing one time-step it is impossible to make any promises of a non exponentially expensive initialization. \\
Furthermore the method of Budinski has no pre-determined quantum circuit for the implementation of the matrix that represents the collision operation. This means that for each timestep a block diagonal matrix $A$ consisting of $32 M \otimes M$ blocks has to be decomposed as a quantum circuit. Both the decomposition calculation and the final decomposed circuit might be exponentially expensive in the number of qubits that the block is applied to, which is another point where any potential quantum advantage might get lost even when performing only one timestep. 


The Carleman linearization function presented in \cite{Sanavio2023} also allows for solving the full Boltzmann equation on a quantum computer. Here first the so-called Carleman linearization is used to linearize the Boltzmann equation. The amount of Carleman variables required is of order $\mathcal{O} \left ( N^k Q^k \right )$, where $k$ is the truncation order \cite{Sanavio2023}. The challenge in this technique comes from then applying the linear operator that evolves the system per time step to the quantum state. This operator then consists of a $\mathcal{O} \left ( N^k Q^k \right )$ size matrix with a non-trivial decomposition, meaning a qiskit decomposition is required. Both the algorithm to decompose the unitary and the resulting circuit are potentially exponentially expensive. In the paper itself they also describe that performing this method for multiple time steps is currently too expensive due to the formidable depth of the resulting circuits.

The space-time encoding has two main benefits compared to the other methods. First of all, all the algorithmic procedures are naturally expressed in quantum gates. There is no need to decompose very large matrices into quantum gates, which as expressed before is exponentially expensive in both computation and potential outcome. The complexity of our streaming steps are $\mathcal{O} \left ( \left ( N_t - t \right )^2  \right )$ with a depth of only $\log_2 \left (N_t -t \right)$ where $N_t$ is capped by $N_g$ (where $N_g$ expresses the maximum number of grid points in a dimension). Furthermore the complexity of the collision step for the D2Q4 case is $\mathcal{O} \left ( \left ( N_t - t \right )^2  \right )$ with a depth of $9$.
